\section{Introduction}
\label{sec:introduction}

Large Language Models (LLMs) based on the Transformer architecture have achieved state-of-the-art performance across automated reasoning, software generation, and conversational interaction. However, serving autoregressive LLMs at low latency and high throughput presents severe computational and memory infrastructure challenges. 

Autoregressive inference inherently exhibits a structural execution duality:
\begin{enumerate}
    \item \textbf{Prefill Phase}: Processes input prompt tokens in parallel. Self-attention computation exhibits quadratic complexity $O(N^2)$ with prompt length, although observed end-to-end prefill latency depends on the full model linear projections, feed-forward layers, and attention kernel implementation.
    \item \textbf{Decode Phase}: Generates tokens sequentially one by one, operating as a memory-bandwidth-dominated workload during single-token autoregressive generation steps.
\end{enumerate}

To alleviate these computational bottlenecks, recent systems research has introduced diverse serving optimizations, including Paged Key-Value (KV) memory allocation~\cite{kwon2023pagedattention}, continuous request batching~\cite{yu2022orca}, radix-tree prefix caching~\cite{zheng2024sglang}, post-training weight and activation quantization~\cite{dettmers2022llmint8,xiao2023smoothquant}, speculative decoding~\cite{chen2023speculative,leviathan2023fast}, and intra-node tensor parallelism~\cite{shoeybi2019megatron}.

Despite these algorithmic advances, existing inference frameworks typically integrate these optimizations as monolithic, highly coupled sub-systems inside complex C++/CUDA serving engines (e.g., vLLM, TensorRT-LLM, SGLang). Consequently, isolating the individual performance contributions, empirical regime boundaries, memory footprint overheads, and non-monotonic combinatorial interaction effects remains an open challenge for system designers.

To address this gap, we introduce \textbf{MicroGen}\footnote{Our open-source framework and benchmark reproducibility package are available at \url{https://github.com/Omdeepb69/MicroGen}}, a modular, hardware-aware execution substrate and empirical research framework designed specifically to isolate, benchmark, and characterize memory-latency-throughput trade-offs in LLM inference optimizations under controlled hardware, workload, and concurrency regimes. Rather than serving as a production engine competing directly with optimized C++ engines, \textsf{MicroGen} provides a clean, protocol-oriented experimental testbed for systematic micro-ablation.

\subsection{Core Research Questions}
\label{sec:rqs}

Our empirical study is guided by four core Research Questions (RQs):
\begin{itemize}
    \item \textbf{RQ1 (Optimization Efficacy Regimes)}: Under what specific sequence lengths, context depths, and prompt structures do individual inference optimizations (Paged KV, Prefix Caching, INT8 Quantization, Speculative Decoding) yield positive latency speedups versus execution overheads?
    \item \textbf{RQ2 (Memory/Throughput Trade-offs)}: What are the exact memory footprint trade-offs (allocated vs. reserved VRAM, quantization lifecycle overhead, KV cache block fragmentation) across model parameter scales?
    \item \textbf{RQ3 (Serving Under Realistic Load)}: How do continuous batching, paged memory allocation, and prefix caching perform under dynamic concurrent request streams ($B \in [1..16]$)?
    \item \textbf{RQ4 (Optimization Regime Generalization)}: Can empirical measurements characterize global workload regime maps, and do these optimization relationships generalize across distinct model parameter scales and architectures?
\end{itemize}

\subsection{Key Contributions}
\label{sec:contributions}

This paper makes the following key contributions:
\begin{enumerate}
    \item \textbf{Modular Experimental Substrate}: We design and implement \textsf{MicroGen}, a Python/PyTorch inference research substrate that cleanly decouples execution backends (\textsf{InferenceBackend}), hardware allocators (\textsf{Device}), KV cache managers (\textsf{KVCache}), continuous batching schedulers (\textsf{Scheduler}), and draft-target verification loops (\textsf{SpeculativeEngine}).
    \item \textbf{Rigorous Benchmark \& Verification Protocol}: We establish a 5-step hardware verification protocol enforcing $N=30$ statistically significant trials, CUDA host-device synchronization, explicit garbage collection, and strictly monotonic microsecond-level latency instrumentation using high-resolution monotonic timers.
    \item \textbf{Empirical Trade-off \& Regime Characterization}: Based on extensive evaluation on NVIDIA Tesla T4 GPUs across open models (\texttt{sshleifer/tiny-gpt2} and \texttt{gpt2}), we empirically demonstrate:
    \begin{itemize}
        \item \textbf{Prefix Caching}: Achieves up to \textbf{3.91$\times$ prefill TTFT speedup} (reducing prefill TTFT from $25.8 \pm 1.2\text{ ms}$ to $6.6 \pm 0.3\text{ ms}$, $p < 0.001$) under 100\% prompt overlap at $L_{\text{prompt}}=1024$ tokens.
        \item \textbf{Paged KV Allocation}: Eliminates external memory fragmentation ($0.0\%$ fragmentation) under dynamic allocation workloads.
        \item \textbf{Non-Monotonic Optimization Composition}: We demonstrate that optimization composition is non-monotonic: combining optimizations (+INT8, +Paged KV, +Prefix Cache) yields a statistically significant throughput regression ($493.6 \pm 8.6\text{ tok/s}$, $0.96\times$ baseline, $p < 0.001$), while full serving stack concurrency drops throughput to $0.76\times$ baseline ($391.8 \pm 6.2\text{ tok/s}$, $p < 0.001$).
    \end{itemize}
    \item \textbf{Open-Source Research Package}: We release the complete \textsf{MicroGen} codebase, experimental manifests, export pipelines, and reproducibility package to facilitate transparent research in LLM inference systems.
\end{enumerate}

\subsection{Paper Organization}
\label{sec:organization}

The remainder of this paper is organized as follows: Section~\ref{sec:architecture} presents the modular architecture of \textsf{MicroGen}. Section~\ref{sec:optimizations} details the mathematical mechanics of our implemented inference optimizations. Section~\ref{sec:methodology} describes our experimental methodology and verification protocol. Section~\ref{sec:evaluation} presents empirical results answering RQ1–RQ4. Section~\ref{sec:related_work} discusses related work, Section~\ref{sec:discussion} analyzes threats to validity and design guidelines, and Section~\ref{sec:conclusion} concludes the paper.
